\section{Battery Theorem Validation}\label{sec:theorem_validation_longform}\label{app:integrated-theorem-surface-register}

The earlier sections state the main theorem surfaces in compact paper form.
This long-form block mirrors the thesis theorem ladder more faithfully by
pulling the full battery-domain theorem chapters into the primary manuscript.
The goal is not to change the bounded claim discipline of the paper, but to
make the theorem scaffolding, assumptions, and proof narratives visible inside
the same canonical `paper.pdf` build rather than only in the deprecated
reference thesis build.

\PaperImportStructuredFile{../chapters/ch15_assumptions_notation_proof_discipline.tex}
\PaperImportStructuredFile{../chapters/ch16_battery_theorem_oasg_existence.tex}
\PaperImportStructuredFile{../chapters/ch17_battery_theorem_safety_preservation.tex}
\PaperImportStructuredFile{../chapters/ch18_orius_core_bound_battery.tex}
\PaperImportStructuredFile{../chapters/ch19_no_free_safety_battery.tex}
\PaperImportStructuredFile{../chapters/ch20_temporal_behavioral_extensions.tex}
\PaperImportStructuredFile{longform/ch21_conditional_coverage_theorem.tex}
\PaperImportStructuredFile{longform/ch22_byzantine_robustness.tex}
\PaperImportStructuredFile{longform/ch23_committee_response.tex}
\PaperImportStructuredFile{longform/ch24_universal_portability.tex}
\PaperImportStructuredFile{longform/ch25_deep_theoretical_guarantees.tex}

\paragraph{From battery ladder to ORIUS program.}
The theorem ladder above is Paper~1's contribution: a locked battery proof
witness anchored by T1--T8.  Papers~2--6 extend specific rungs of that ladder
into the five ORIUS program layers:
\begin{itemize}[leftmargin=*,topsep=2pt,itemsep=1pt]
  \item \textbf{Paper~2 (temporal validity):} takes the certificate horizon
        theorem (T5) and the expiration bound (T6) and turns them into a
        runtime validity-horizon engine
        (\texttt{dc3s/half\_life.py}, \texttt{dc3s/reachability.py}).
  \item \textbf{Paper~3 (graceful fallback):} takes the fallback existence
        theorem (T7) and the zero-dispatch safe-hold and builds a constructive
        four-policy degradation planner (\texttt{dc3s/graceful.py}).
  \item \textbf{Paper~4 (benchmark discipline):} formalizes the evaluation
        contract that the battery witness depends on into a replayable,
        schema-governed ORIUS-Bench (\texttt{orius\_bench/}).
  \item \textbf{Paper~5 (composition):} shows that the local certificate
        guarantee from the battery ladder does \emph{not} compose under shared
        feeder constraints, and provides centralized and distributed repair
        protocols (\texttt{multi\_agent/}).
  \item \textbf{Paper~6 (runtime governance):} wraps the full certificate
        lifecycle in a CertOS runtime with three invariants (INV-1, INV-2,
        INV-3) and an auditable hash chain (\texttt{certos/}).
\end{itemize}
Each extension stays bounded: Papers~2 and~3 are theorem-backed within the
battery domain; Paper~4 is a benchmark-contract contribution; Paper~5 is a
bounded composition counterexample; Paper~6 is a bounded runtime-governance
layer.  They do not upgrade bounded evidence into theorem closure.

